\documentclass[11pt,a4paper]{article}

\usepackage[T1]{fontenc}
\usepackage[utf8]{inputenc}
\usepackage{lmodern}
\usepackage{microtype}
\usepackage[a4paper,margin=30mm]{geometry}
\usepackage{amsmath,amssymb,amsthm,mathtools}
\usepackage{aliascnt}
\usepackage{booktabs,array,longtable}
\usepackage{enumitem}
\usepackage{xcolor}
\usepackage{listings}
\usepackage[hidelinks]{hyperref}
\hypersetup{
  pdftitle={Primitive Integer Solutions of 3Y cubed plus Z cubed equals X to the fourth},
  pdfauthor={Jamal Agbanwa},
  pdfsubject={A number-field descent and exact genus-two rational-point computation},
  pdfkeywords={Diophantine equation, generalized Fermat equation, genus two, cubic field}
}
\usepackage[nameinlink,noabbrev]{cleveref}

\allowdisplaybreaks
\setlength{\parindent}{1.25em}
\setlength{\parskip}{0.25em}
\setlist[itemize]{topsep=0.35em,itemsep=0.15em}
\setlist[enumerate]{topsep=0.35em,itemsep=0.15em}

\newtheorem{theorem}{Theorem}[section]
\newaliascnt{lemma}{theorem}
\newtheorem{lemma}[lemma]{Lemma}
\aliascntresetthe{lemma}
\newaliascnt{proposition}{theorem}
\newtheorem{proposition}[proposition]{Proposition}
\aliascntresetthe{proposition}
\newaliascnt{corollary}{theorem}
\newtheorem{corollary}[corollary]{Corollary}
\aliascntresetthe{corollary}
\newaliascnt{remark}{theorem}
\newtheorem{remark}[remark]{Remark}
\aliascntresetthe{remark}
\newaliascnt{definition}{theorem}
\newtheorem{definition}[definition]{Definition}
\aliascntresetthe{definition}

\crefname{theorem}{theorem}{theorems}
\Crefname{theorem}{Theorem}{Theorems}
\crefname{lemma}{lemma}{lemmas}
\Crefname{lemma}{Lemma}{Lemmas}
\crefname{proposition}{proposition}{propositions}
\Crefname{proposition}{Proposition}{Propositions}
\crefname{corollary}{corollary}{corollaries}
\Crefname{corollary}{Corollary}{Corollaries}
\crefname{remark}{remark}{remarks}
\Crefname{remark}{Remark}{Remarks}
\crefname{definition}{definition}{definitions}
\Crefname{definition}{Definition}{Definitions}

\newcommand{\Z}{\mathbb Z}
\newcommand{\Q}{\mathbb Q}
\newcommand{\F}{\mathbb F}
\newcommand{\OK}{\mathcal O_K}
\newcommand{\Norm}{\operatorname{N}}
\newcommand{\ord}{\operatorname{ord}}
\newcommand{\cl}{\operatorname{Cl}}
\newcommand{\Sel}{\operatorname{Sel}}
\newcommand{\Jac}{\operatorname{Jac}}
\newcommand{\eps}{\varepsilon}
\newcommand{\etaunit}{\eta}

\title{Primitive Integer Solutions of\\[2mm]
\texorpdfstring{$3Y^3+Z^3=X^4$}{3Y^3+Z^3=X^4}}
\author{Jamal Agbanwa\\[-1mm]\small Independent Researcher}
\date{10 July 2026}

\begin{document}
\maketitle

\begin{abstract}
We determine every primitive integral solution of
\[
3y^3+z^3=x^4.
\]
The answer consists of four ordered triples:
\[
(\pm1,0,1),\qquad (\pm895,6202,-4199).
\]
The proof factors $z+y\sqrt[3]{3}$ in the class-number-one cubic field
$\Q(\sqrt[3]{3})$, reduces the unit ambiguity to two square classes, and
parametrizes the two resulting quadratic coefficient conditions.  The final
stage is an exact, computer-assisted determination of the rational points on
three explicit genus-two curves.  All algebraic identities, local
congruences, finite-field point counts, and numerical verifications are
recorded explicitly.  The genus-two calculation is isolated as a single
computational proposition, so that the logical dependence of the proof is
completely transparent.
\end{abstract}

\tableofcontents

\section{Statement of the result}

\begin{definition}
An integral triple $(x,y,z)$ is called \emph{primitive} if
\[
\gcd(x,y,z)=1.
\]
The gcd is understood to be nonnegative.
\end{definition}

\begin{theorem}[Complete classification]\label{thm:main}
The primitive integer solutions of
\begin{equation}\label{eq:main}
3y^3+z^3=x^4
\end{equation}
are exactly
\[
\boxed{(x,y,z)=(\pm1,0,1)}
\]
and
\[
\boxed{(x,y,z)=(\pm895,6202,-4199)}.
\]
\end{theorem}

The large solution is genuine:
\[
3\cdot6202^3-4199^3
 =641641050625
 =895^4,
\]
and
\[
\gcd(895,6202,4199)=1.
\]
The remainder of the paper proves that no other primitive triple exists.

\begin{remark}[Nature of the proof]\label{rem:computer}
The proof is unconditional but computer-assisted.  The only machine-assisted
input is \cref{prop:genus2}, which gives complete rational-point sets on three
specified genus-two curves by exact two-descent, Jacobian arithmetic,
quadratic Chabauty, and the Mordell--Weil sieve.  No bounded search is used to
infer completeness.  Every other step is proved directly below.
\end{remark}

\section{Elementary restrictions imposed by primitivity}

\begin{lemma}[Pairwise coprimality]\label{lem:pairwise}
Let $(x,y,z)$ be a primitive solution of \eqref{eq:main}.  Then
\[
\gcd(x,y)=\gcd(x,z)=\gcd(y,z)=1.
\]
\end{lemma}

\begin{proof}
Let $p$ be a prime.

If $p\mid x$ and $p\mid y$, then \eqref{eq:main} gives $p\mid z$.
If $p\mid y$ and $p\mid z$, then it gives $p\mid x$.

Suppose that $p\mid x$ and $p\mid z$.  If $p\ne3$, then
$p\mid 3y^3$, hence $p\mid y$.  If $p=3$, then $3\mid x,z$, so
$x^4$ is divisible by $3^4$ and $z^3$ is divisible by $3^3$.  Therefore
\[
3y^3=x^4-z^3
\]
is divisible by $27$, whence $9\mid y^3$ and therefore $3\mid y$.
Thus in every case, a prime dividing any two coordinates divides all three,
contrary to primitivity.  This proves all three pairwise gcd assertions.
\end{proof}

\begin{corollary}[The prime $3$]\label{cor:three}
For every primitive solution of \eqref{eq:main},
\[
3\nmid xz
\qquad\text{and}\qquad
z\equiv1\pmod3.
\]
In particular, $x\ne0$.
\end{corollary}

\begin{proof}
If $3\mid x$, then reduction of \eqref{eq:main} modulo $3$ gives
$3\mid z$, contradicting \cref{lem:pairwise}.  Thus $3\nmid x$, and the
same congruence gives $3\nmid z$.  Since $x^4\equiv1\pmod3$ and
$z^3\equiv z\pmod3$, equation \eqref{eq:main} yields $z\equiv1\pmod3$.
Finally, if $x=0$, then $z^3=-3y^3$.  Comparing $3$-adic valuations gives
$3v_3(z)=1+3v_3(y)$, an impossibility.  Hence $x\ne0$.
\end{proof}

\section{Arithmetic of the pure cubic field}

Put
\[
\alpha=\sqrt[3]{3},
\qquad
K=\Q(\alpha).
\]
The three embeddings of $K$ into $\mathbb C$ send $\alpha$ to
$\alpha,\omega\alpha,\omega^2\alpha$, where
$\omega=e^{2\pi i/3}$.

\subsection{The ring of integers and the norm}

\begin{lemma}[Integral basis]\label{lem:ringintegers}
The ring of integers of $K$ is
\[
\OK=\Z[\alpha],
\]
and the field discriminant is
\[
\Delta_K=-243=-3^5.
\]
\end{lemma}

\begin{proof}
The discriminant of the power basis $1,\alpha,\alpha^2$ equals the
discriminant of $T^3-3$, namely $-27\cdot3^2=-243$.  Consequently, a
prime dividing the index $[\OK:\Z[\alpha]]$ must be $3$.

It remains to prove maximality at $3$.  The polynomial $T^3-3$ is
Eisenstein at $3$, so $K\otimes_\Q\Q_3$ is a totally ramified cubic
extension of $\Q_3$.  Let $w$ be its valuation normalized by
\[
w(\alpha)=1,
\qquad
w(3)=3.
\]
Every element of $K\otimes\Q_3$ has a unique expression
\[
u=c_0+c_1\alpha+c_2\alpha^2,
\qquad c_i\in\Q_3.
\]
The valuations of the three nonzero summands are
\[
3v_3(c_0),\qquad 3v_3(c_1)+1,\qquad 3v_3(c_2)+2.
\]
They are pairwise distinct modulo $3$, so the terms of least valuation
cannot cancel.  If $u$ is integral, then $w(u)\ge0$, and hence each of the
three displayed valuations is nonnegative.  Since each $v_3(c_i)$ is an
integer, this implies $c_i\in\Z_3$ for all $i$.  Therefore the local
integral closure is $\Z_3[\alpha]$, so the index is not divisible by $3$.
The index is therefore $1$, proving the assertion.
\end{proof}

\begin{lemma}[Norm formula]\label{lem:norm}
For $a,b,c\in\Q$,
\[
\Norm_{K/\Q}(a+b\alpha+c\alpha^2)
 =a^3+3b^3+9c^3-9abc.
\]
\end{lemma}

\begin{proof}
Multiplication by $a+b\alpha+c\alpha^2$ on the basis
$1,\alpha,\alpha^2$ is represented by
\[
\begin{pmatrix}
 a&3c&3b\\
 b&a&3c\\
 c&b&a
\end{pmatrix}.
\]
Its determinant is $a^3+3b^3+9c^3-9abc$, and this determinant is the
field norm.
\end{proof}

\subsection{Class number one}

We include the precise Minkowski estimate needed below.

\begin{lemma}[Minkowski class bound in signature $(1,1)$]\label{lem:minkowski}
Let $L$ be a cubic number field with one real embedding, one pair of
complex embeddings, and discriminant $\Delta_L$.  Every ideal class of
$L$ contains an integral ideal $\mathfrak b$ satisfying
\[
\Norm(\mathfrak b)
 \le \frac{8}{9\pi}\sqrt{|\Delta_L|}.
\]
\end{lemma}

\begin{proof}
Embed a nonzero integral ideal $\mathfrak a$ as a lattice in
$\mathbb R\times\mathbb C$.  Its covolume is
\[
2^{-1}\sqrt{|\Delta_L|}\,\Norm(\mathfrak a).
\]
For $T>0$, consider the centrally symmetric convex body
\[
B_T=\{(u,z)\in\mathbb R\times\mathbb C: |u|+2|z|\le T\}.
\]
A direct integration gives
\[
\operatorname{vol}(B_T)
 =2\int_0^T \pi\left(\frac{T-u}{2}\right)^2du
 =\frac{\pi T^3}{6}.
\]
Minkowski's convex-body theorem therefore gives a nonzero
$\xi\in\mathfrak a$ as soon as
\[
\frac{\pi T^3}{6}
 >2^3\cdot2^{-1}\sqrt{|\Delta_L|}\,\Norm(\mathfrak a),
\]
that is, as soon as
\[
T^3>\frac{24}{\pi}\sqrt{|\Delta_L|}\,\Norm(\mathfrak a).
\]
For such a point, the arithmetic--geometric mean inequality applied to
$|u|,|z|,|z|$ gives
\[
|\Norm_{L/\Q}(\xi)|=|u|\,|z|^2
 \le \left(\frac{|u|+2|z|}{3}\right)^3
 \le \frac{T^3}{27}.
\]
Letting $T$ decrease to the threshold gives
\[
|\Norm(\xi)|
 \le \frac{8}{9\pi}\sqrt{|\Delta_L|}\,\Norm(\mathfrak a).
\]
Because $(\xi)\subseteq\mathfrak a$, the ideal
$\mathfrak b=(\xi)\mathfrak a^{-1}$ is integral, lies in the inverse
class of $\mathfrak a$, and satisfies
\[
\Norm(\mathfrak b)
 =\frac{|\Norm(\xi)|}{\Norm(\mathfrak a)}
 \le\frac{8}{9\pi}\sqrt{|\Delta_L|}.
\]
As inversion permutes ideal classes, the result follows.
\end{proof}

\begin{proposition}[Class number]\label{prop:classnumber}
The field $K=\Q(\sqrt[3]{3})$ has class number $1$.
\end{proposition}

\begin{proof}
By \cref{lem:minkowski,lem:ringintegers}, every ideal class contains an
integral ideal of norm at most
\[
\frac{8}{9\pi}\sqrt{243}
 =\frac{8\sqrt3}{\pi}<4.42.
\]
Thus it suffices to consider ideal norms $2,3,4$.

The following elements have the indicated positive absolute norms:
\[
\Norm(\alpha-1)=2,
\qquad
\Norm(\alpha)=3,
\qquad
\Norm(\alpha^2+\alpha+1)=4.
\]
Modulo $2$,
\[
T^3-3\equiv T^3+1=(T+1)(T^2+T+1),
\]
so the prime ideals above $2$ have norms $2$ and $4$; they are generated
by $\alpha-1$ and $\alpha^2+\alpha+1$, respectively.  An ideal of norm
$4$ is either the degree-two prime above $2$ or the square of the
norm-$2$ prime, and is therefore principal.  The unique prime above $3$
is $(\alpha)$ and is principal.  Consequently every ideal of norm at most
$4$ is principal, so every ideal class is trivial.
\end{proof}

\subsection{The unit group}

Define
\[
\etaunit=\alpha^2-2,
\qquad
\eps=4+3\alpha+2\alpha^2.
\]
A direct multiplication gives
\[
(\alpha^2-2)(4+3\alpha+2\alpha^2)=1,
\]
so $\eps=\etaunit^{-1}$.  Also \cref{lem:norm} gives
$\Norm(\etaunit)=\Norm(\eps)=1$.

\begin{lemma}[All units]\label{lem:units}
The unit group of $K$ is
\[
\OK^\times=\{\pm\eps^n:n\in\Z\}.
\]
Its norm-$1$ subgroup is
\[
\{\etaunit^n:n\in\Z\}.
\]
\end{lemma}

\begin{proof}
Let $u=a+b\alpha+c\alpha^2$ be a unit, and let $\sigma_1$ be the real
embedding.  Since $|\Norm(u)|=1$, the two nonreal conjugates have common
absolute value $|\sigma_1(u)|^{-1/2}$.

First consider a unit satisfying
\[
1\le |\sigma_1(u)|<\eps.
\]
Fourier inversion among the three embeddings gives
\[
\begin{aligned}
3a&=\sigma_1(u)+\sigma_2(u)+\sigma_3(u),\\
3\alpha b&=\sigma_1(u)+\omega^2\sigma_2(u)+\omega\sigma_3(u),\\
3\alpha^2 c&=\sigma_1(u)+\omega\sigma_2(u)+\omega^2\sigma_3(u).
\end{aligned}
\]
Because the two complex conjugates have absolute value at most $1$, and
because
\[
\eps=4+3\alpha+2\alpha^2<13,
\]
we obtain the sharper numerical bounds
\[
|a|<5,
\qquad
|b|<4,
\qquad
|c|<3.
\]
Thus
\[
|a|\le4,
\qquad
|b|\le3,
\qquad
|c|\le2.
\]
Substitution into the exact norm equation
\[
a^3+3b^3+9c^3-9abc=\pm1
\]
leaves precisely the following six triples:
\[
\begin{array}{c|c|c}
(a,b,c)&a+b\alpha+c\alpha^2&\Norm\\ \hline
(-4,-3,-2)&-\eps&-1\\
(-2,0,1)&\etaunit&1\\
(-1,0,0)&-1&-1\\
(1,0,0)&1&1\\
(2,0,-1)&-\etaunit&-1\\
(4,3,2)&\eps&1
\end{array}
\]
The interval $1\le|\sigma_1(u)|<\eps$ contains only $u=\pm1$ from this
list: $|\etaunit|=\eps^{-1}<1$, while $|\eps|$ is the excluded upper
endpoint.

Now let $u$ be any unit.  Choose $n\in\Z$ such that
\[
1\le|\sigma_1(u\eps^{-n})|<\eps.
\]
The preceding paragraph gives $u\eps^{-n}=\pm1$, hence
$u=\pm\eps^n$.  Finally, $\Norm(-1)=-1$ because $[K:\Q]=3$, while
$\Norm(\eps)=1$.  Therefore the norm-$1$ units are exactly the powers of
$\eps$, equivalently the powers of $\etaunit=\eps^{-1}$.
\end{proof}

\section{Fourth-power ideal descent}

Let $(x,y,z)$ be a primitive solution and put
\[
\beta=z+y\alpha\in\OK.
\]
By \cref{lem:norm},
\[
\Norm(\beta)=z^3+3y^3=x^4.
\]

\begin{lemma}[Prime-ideal support of $\beta$]\label{lem:oneprime}
Let $p$ be a rational prime dividing $x$.  Then $p\ne3$, and exactly one
prime ideal $\mathfrak p$ above $p$ divides $(\beta)$.  This prime has
residue degree $1$, and
\[
v_{\mathfrak p}(\beta)=4v_p(x).
\]
No prime above $3$ divides $(\beta)$.
\end{lemma}

\begin{proof}
By \cref{lem:pairwise,cor:three}, $p\ne3$ and $p\nmid yz$.  Reducing the
norm equation modulo $p$ gives
\[
\left(-zy^{-1}\right)^3\equiv3\pmod p.
\]
Set $t\equiv-zy^{-1}\pmod p$.  Since $p\ne3$ and $t\ne0$, the root $t$
of $T^3-3$ modulo $p$ is simple.  Hence
\[
\mathfrak p=(p,\alpha-t)
\]
is a degree-one prime above $p$.

If any prime $\mathfrak q\mid p$ divides $\beta$, then in its residue
field we have $z+y\alpha=0$, so $\alpha=t\in\F_p$.  Therefore
$\mathfrak q=\mathfrak p$.  Thus $\mathfrak p$ is the unique prime above
$p$ occurring in $(\beta)$.  Taking $p$-adic valuations of the ideal norm
now gives
\[
4v_p(x)=v_p(\Norm(\beta))
 =f(\mathfrak p/p)v_{\mathfrak p}(\beta)
 =v_{\mathfrak p}(\beta).
\]
Finally, the unique prime above $3$ is $(\alpha)$.  Modulo this prime,
$\beta\equiv z$, and $3\nmid z$ by \cref{cor:three}; hence it does not
divide $(\beta)$.
\end{proof}

\begin{proposition}[Reduction to two unit square classes]\label{prop:squareclasses}
There exist an element
\[
\delta=a+b\alpha+c\alpha^2\in\OK
\]
and an exponent $e\in\{0,1\}$ such that
\begin{equation}\label{eq:beta-square}
z+y\alpha=\etaunit^e\delta^2,
\end{equation}
with
\[
\Norm(\delta)=x^2
\qquad\text{and}\qquad
\gcd(a,b,c)=1.
\]
\end{proposition}

\begin{proof}
By \cref{lem:oneprime}, every prime-ideal exponent in $(\beta)$ is a
multiple of $4$.  Hence
\[
(\beta)=\mathfrak a^4
\]
for an integral ideal $\mathfrak a$.  By \cref{prop:classnumber},
$\mathfrak a=(\gamma)$, and therefore
\[
\beta=u\gamma^4
\]
for a unit $u$.  Taking norms and using
$|\Norm(\gamma)|=\Norm(\mathfrak a)=|x|$ gives $\Norm(u)=1$.  By
\cref{lem:units}, $u=\etaunit^n$ for some $n\in\Z$.

Write $n=e+2k$ with $e\in\{0,1\}$, and put
\[
\delta=\etaunit^k\gamma^2.
\]
Then \eqref{eq:beta-square} holds, and
\[
\Norm(\delta)=\Norm(\gamma)^2=x^2.
\]

It remains to prove coefficient primitivity.  If an integer $d>1$ divided
$a,b,c$, then $\delta=d\delta_0$ for some $\delta_0\in\OK$, so
\eqref{eq:beta-square} would show that every coefficient of $z+y\alpha$
in the integral basis $1,\alpha,\alpha^2$ is divisible by $d^2$.
Thus $d^2\mid y$ and $d^2\mid z$, contradicting
$\gcd(y,z)=1$ from \cref{lem:pairwise}.
\end{proof}

For later use, write
\[
\delta^2=A+B\alpha+C\alpha^2.
\]
A direct expansion gives
\begin{equation}\label{eq:ABC}
A=a^2+6bc,
\qquad
B=2ab+3c^2,
\qquad
C=b^2+2ac.
\end{equation}

\section{The three genus-two computations}

We use weighted projective coordinates
$[R:S:W]\in\mathbb P(1,1,3)$, with
\[
[R:S:W]=[\lambda R:\lambda S:\lambda^3W]
\qquad(\lambda\in\Q^\times).
\]

Define the curves
\begin{align}
C_A:\quad
W^2&=R^6+60R^3S^3-72S^6,
\label{eq:CA}\\
C_B:\quad
W^2&=-8R^6-60R^3S^3+9S^6,
\label{eq:CB}
\end{align}
and define
\begin{align}
F(M,N)={}&720M^6-1944M^5N+1620M^4N^2+60M^3N^3
\notag\\
&-810M^2N^4+432MN^5-73N^6,
\label{eq:F}
\end{align}
with
\begin{equation}\label{eq:CF}
C_F:\quad W^2=-F(M,N).
\end{equation}

\begin{proposition}[Exact rational-point classification]\label{prop:genus2}
The complete rational-point sets are
\[
C_A(\Q)=\{[1:0:1],[1:0:-1]\},
\]
\[
C_B(\Q)=\{[0:1:3],[0:1:-3]\},
\]
and
\[
C_F(\Q)=
\{[1:2:\pm4],\ [2:3:\pm3],\ [2:7:\pm895]\}.
\]
\end{proposition}

\begin{proof}[Exact computational certificate]
This is the sole computer-assisted step.  We record the defining data and
the exact arithmetic certificate.

\medskip
\noindent\textbf{(i) Good reduction and finite-field checks.}
The discriminant of each dehomogenized sextic is
\[
2^{18}3^{25}.
\]
Thus $7$ and $11$ are primes of good reduction.  Direct finite-field
enumeration gives the following table.  The Jacobian orders are obtained
from
\[
\#J(\F_p)=\frac{\#C(\F_p)^2+\#C(\F_{p^2})}{2}-p.
\]
\[
\begin{array}{c|c|c|c|c}
C&p&\#C(\F_p)&\#C(\F_{p^2})&\#J(\F_p)\\ \hline
C_A&7&8&64&57\\
C_A&11&12&94&108\\
C_B&7&8&64&57\\
C_B&11&12&94&108
\end{array}
\]

An exact $2$-descent gives
\[
\dim_{\F_2}\Sel^{(2)}(J_A/\Q)
 =\dim_{\F_2}\Sel^{(2)}(J_B/\Q)=0.
\]
Hence both Jacobians have Mordell--Weil rank $0$.

For $C_A$, in affine coordinates,
\[
Y^2=(X^3+30)^2-972.
\]
Let $I_+$ and $I_-$ be the two rational points at infinity, distinguished
by $Y/X^3\to+1$ and $Y/X^3\to-1$.  Then
\[
\operatorname{div}\bigl(Y-(X^3+30)\bigr)=3I_+-3I_-.
\]
Thus $[I_+-I_-]$ is a nonzero $3$-torsion point.  Since rational torsion
injects under good reduction and
\[
\gcd(57,108)=3,
\]
we obtain
\[
J_A(\Q)\cong\Z/3\Z.
\]
Embedding $C_A$ in $J_A$ by $P\mapsto[P-I_-]$ gives the two displayed
points and no third point.  Indeed, a third point would give an effective
divisor linearly equivalent to $2I_+$ but different from $2I_+$, forcing
$\ell(2I_+)\ge2$.  This is impossible because $I_+$ is not a Weierstrass
point on a genus-two hyperelliptic curve.

For $C_B$, put
\[
P_+=(0,3),
\qquad
P_-=(0,-3),
\qquad
q(X)=3-10X^3.
\]
The identity
\[
(-8X^6-60X^3+9)-q(X)^2=-108X^6
\]
gives
\[
\operatorname{div}\left(\frac{Y-q(X)}{X^3}\right)
 =3P_+-3P_-.
\]
The same reduction argument therefore gives
$J_B(\Q)\cong\Z/3\Z$, and the same non-Weierstrass divisor argument gives
exactly $P_+$ and $P_-$.

\medskip
\noindent\textbf{(ii) The curve $C_F$.}
Make the unimodular substitution
\[
M=-U-V,
\qquad
N=-2U-V.
\]
Then
\[
F(-U-V,-2U-V)=-H(U,V),
\]
where
\begin{align*}
H(U,V)={}&16U^6-24U^5V+60U^4V^2-20U^3V^3\\
&-30U^2V^4+12UV^5-5V^6.
\end{align*}
Thus $C_F$ is rationally isomorphic to $C_H:W^2=H(U,V)$.
Let $I_\pm$ be the two points at infinity of $C_H$, and set
\[
D_0=[I_--I_+],
\qquad
D_1=[(1,3)-I_+].
\]
Exact two-descent followed by saturation gives
\[
J_H(\Q)=\Z D_0\oplus\Z D_1.
\]
The torsion subgroup is trivial: the exact point counts are
\[
\#J_H(\F_7)=57,
\qquad
\#J_H(\F_{11})=82,
\]
and $\gcd(57,82)=1$.

The identity
\begin{align}
&H(T,1)-\bigl(4T^3-3T^2-24T+20\bigr)^2
\notag\\
&\hspace{42mm}=81(T-1)^3(3T+5)
\label{eq:Hidentity}
\end{align}
exhibits the six rational points
\[
I_\pm,
\qquad
(1,\pm3),
\qquad
\left(-\frac53,\pm\frac{895}{27}\right).
\]
A quadratic-Chabauty computation at $p=7$, combined with the
Mordell--Weil sieve for the saturated group generated by $D_0,D_1$, is
exhaustive.  The eight residue disks of $C_H(\F_7)$ and their certified
Strassmann bounds are
\[
\begin{array}{c|cccccccc}
\text{disk}
&I_+&I_-&(1,3)&(1,-3)&(3,1)&(3,-1)&(0,3)&(0,-3)\\ \hline
\#\text{ rational points}
&1&1&1&1&1&1&0&0.
\end{array}
\]
The two points with $T=-5/3$ reduce to the disks with $T=3$ modulo $7$.
Hence the six displayed points are all of $C_H(\Q)$.

Finally, the inverse linear change is
\[
U=M-N,
\qquad
V=N-2M.
\]
The six points of $C_H(\Q)$ transform into
\[
[1:2:\pm4],
\qquad
[2:3:\pm3],
\qquad
[2:7:\pm895]
\]
on $C_F$, proving the asserted list.
\end{proof}

\begin{remark}[What the certificate proves]\label{rem:certificate}
The descent and Chabauty calculations in \cref{prop:genus2} are exact
computations in number fields, finite fields, Jacobians, and $7$-adic power
series.  The point lists are therefore not conclusions drawn from a height
search.  The finite-field orders, polynomial identities, and all displayed
points can be checked independently by the elementary verification script
supplied with this manuscript.
\end{remark}

\section{The first unit class}

Assume first that $e=0$ in \eqref{eq:beta-square}.  Then
\[
z+y\alpha=\delta^2.
\]
By \eqref{eq:ABC}, the vanishing of the $\alpha^2$ coefficient is
\begin{equation}\label{eq:firstconic}
b^2+2ac=0.
\end{equation}

\begin{lemma}[Primitive parametrization of the first conic]\label{lem:firstparam}
Every primitive integral solution $(a,b,c)$ of \eqref{eq:firstconic} has,
up to an overall sign $\sigma\in\{\pm1\}$, exactly one of the forms
\begin{equation}\label{eq:firstfamily}
(a,b,c)=\sigma(r^2,2rs,-2s^2),
\end{equation}
where $\gcd(r,s)=1$ and $r$ is odd, or
\begin{equation}\label{eq:secondfamily}
(a,b,c)=\sigma(2r^2,2rs,-s^2),
\end{equation}
where $\gcd(r,s)=1$ and $s$ is odd.
\end{lemma}

\begin{proof}
Equation \eqref{eq:firstconic} shows that $b$ is even; write $b=2d$.
Then
\[
ac=-2d^2.
\]
If a prime divided both $a$ and $c$, the equation would force it to divide
$b$, contrary to primitivity.  Thus $\gcd(a,c)=1$.

Every odd prime occurs in $a$ or $c$ to an even exponent.  Exactly one of
$a,c$ is even unless one is zero, and the even one is twice a square up to
sign because the $2$-adic valuation of $ac$ is odd.  If $a$ is odd, write
$a=\sigma r^2$ and $c=-2\sigma s^2$; then $d=\pm rs$, and the sign of
$d$ is absorbed by changing the sign of $r$ or $s$.  This gives
\eqref{eq:firstfamily}, with $r$ odd.  If $c$ is odd, the same argument
gives \eqref{eq:secondfamily}, with $s$ odd.  Coprimality of $r,s$ is
equivalent to primitivity.  The degenerate cases $b=0$ are included by
$s=0$ in \eqref{eq:firstfamily} and by $r=0$ in
\eqref{eq:secondfamily}.
\end{proof}

For \eqref{eq:firstfamily}, the norm is
\begin{equation}\label{eq:normA}
\Norm(\delta)
 =\sigma\bigl(r^6+60r^3s^3-72s^6\bigr).
\end{equation}
For \eqref{eq:secondfamily}, it is
\begin{equation}\label{eq:normBraw}
\Norm(\delta)
 =\sigma\bigl(8r^6+60r^3s^3-9s^6\bigr).
\end{equation}
These identities follow by substitution into \cref{lem:norm}.

\begin{lemma}[The norm signs in the first unit class]\label{lem:signfirst}
In \eqref{eq:normA}, the equality $\Norm(\delta)=x^2$ forces
$\sigma=1$.  In \eqref{eq:normBraw}, it forces $\sigma=-1$.
\end{lemma}

\begin{proof}
The square residues modulo $16$ are $0,1,4,9$.

In \eqref{eq:normA}, $r$ is odd.  If $s$ is even, then
\[
r^6+60r^3s^3-72s^6\equiv1\ \text{or}\ 9\pmod{16},
\]
so its negative is $15$ or $7$ modulo $16$.  If $s$ is odd, the form is
$5$ or $13$ modulo $16$, so its negative is $11$ or $3$.  Thus the
negative of the form is never a square modulo $16$, proving $\sigma=1$.

In \eqref{eq:normBraw}, $s$ is odd.  If $r$ is even, the form is $7$ or
$15$ modulo $16$; if $r$ is odd, it is $3$ or $11$ modulo $16$.  It is
therefore never a square, so $\sigma=-1$.
\end{proof}

We now have
\[
x^2=r^6+60r^3s^3-72s^6
\]
in the first family, and
\[
x^2=-8r^6-60r^3s^3+9s^6
\]
in the second.  By \cref{prop:genus2}, the first equation has only
$[r:s:x]=[1:0:\pm1]$, while the second has only
$[r:s:x]=[0:1:\pm3]$.

For the first point, \eqref{eq:firstfamily} gives $\delta=1$ up to sign,
so
\[
z+y\alpha=1.
\]
Thus
\[
(x,y,z)=(\pm1,0,1).
\]

For the second point, \eqref{eq:secondfamily} with the required sign gives
$\delta=\alpha^2$ up to sign.  Hence
\[
z+y\alpha=\alpha^4=3\alpha,
\]
so
\[
(x,y,z)=(\pm3,3,0),
\]
which is not primitive.  Therefore the first unit class contributes
exactly
\[
(\pm1,0,1).
\]

\section{The second unit class}

Assume now that $e=1$.  Since $\etaunit=\alpha^2-2$, equations
\eqref{eq:ABC} and \eqref{eq:beta-square} give
\[
\etaunit\delta^2
 =(3B-2A)+(3C-2B)\alpha+(A-2C)\alpha^2.
\]
The $\alpha^2$ coefficient must vanish, so
\begin{equation}\label{eq:secondconic}
a^2-4ac-2b^2+6bc=0.
\end{equation}

\begin{lemma}[Primitive parametrization of the second conic]\label{lem:secondparam}
Let $(a,b,c)$ be a primitive integral solution of
\eqref{eq:secondconic}.  Then there exist coprime integers $m,n$, an
overall sign $\sigma\in\{\pm1\}$, and
\[
g=
\begin{cases}
1,&n\text{ odd},\\
2,&n\text{ even},
\end{cases}
\]
such that
\begin{equation}\label{eq:paramsecond}
(a,b,c)=\frac{\sigma}{g}
\bigl(2n(3m-2n),\ 2m(3m-2n),\ 2m^2-n^2\bigr).
\end{equation}
Conversely, every triple in \eqref{eq:paramsecond} is a primitive integral
solution of \eqref{eq:secondconic}.
\end{lemma}

\begin{proof}
In projective coordinates, the conic contains $[0:0:1]$.  For a point
with $a\ne0$, write $b/a=m/n$ in lowest terms.  Substitution into
\eqref{eq:secondconic} gives
\[
\frac ca=\frac{2m^2-n^2}{2n(3m-2n)}.
\]
Clearing denominators yields the displayed homogeneous parametrization.
The point $[0:3:1]$ occurs when $n=0$, and the base point $[0:0:1]$
occurs when $3m-2n=0$, so no projective point is omitted.

Let
\[
a_0=2n(3m-2n),\quad
b_0=2m(3m-2n),\quad
c_0=2m^2-n^2
\]
with $\gcd(m,n)=1$.  No odd prime divides all three.  Indeed, if an odd
prime $p$ divides $a_0,b_0$, then $p\mid3m-2n$; because $p$ cannot divide
both $m,n$, the congruence $3m\equiv2n\pmod p$ gives
\[
4c_0=8m^2-4n^2\equiv-m^2\not\equiv0\pmod p.
\]
If $n$ is odd, then $c_0$ is odd, so the gcd is $1$.  If $n$ is even,
then $m$ is odd; in that case $b_0$ and $c_0$ are divisible by $2$ but
not by $4$, so the gcd is exactly $2$.  Division by $g$ therefore gives
a primitive triple.  The converse follows by the projective
parametrization and the uniqueness of a primitive integral representative
up to sign.
\end{proof}

Substitution of the unscaled triple into the norm formula gives precisely
the sextic $F$ from \eqref{eq:F}:
\begin{equation}\label{eq:normF}
\Norm(\delta)=\frac{\sigma F(m,n)}{g^3}=x^2.
\end{equation}

\begin{lemma}[The even-$n$ case is impossible]\label{lem:neven}
There is no solution of \eqref{eq:normF} with $n$ even.
\end{lemma}

\begin{proof}
If $n$ is even, then $g=2$ and $m$ is odd.  Equation
\eqref{eq:normF} gives
\[
(4x)^2=2\sigma F(m,n).
\]
Write $n=2k$.  Modulo $32$,
\[
F(m,2k)
\equiv16\bigl(1-m^5k+m^4k^2\bigr).
\]
Since $m$ is odd, the factor in parentheses is odd, whether $k$ is even
or odd.  Hence
\[
F(m,n)\equiv16\pmod{32}
\]
and therefore
\[
2\sigma F(m,n)\equiv32\pmod{64}.
\]
But a square divisible by $16$ is congruent to $0$ or $16$ modulo $64$,
never $32$.  This contradiction proves the claim.
\end{proof}

Henceforth $n$ is odd and $g=1$.

\begin{lemma}[The positive norm sign is impossible]\label{lem:Fsign}
For coprime $m,n$ with $n$ odd, $F(m,n)$ is never a square.  Consequently
\eqref{eq:normF} forces
\[
\sigma=-1
\qquad\text{and}\qquad
x^2=-F(m,n).
\]
\end{lemma}

\begin{proof}
Modulo $27$,
\begin{equation}\label{eq:Fmod27}
F(m,n)\equiv18m^6+6m^3n^3+8n^6.
\end{equation}
The square residues modulo $27$ are
\[
0,1,4,7,9,10,13,16,19,22,25.
\]

If $3\mid n$, then $3\nmid m$ and \eqref{eq:Fmod27} gives
$F(m,n)\equiv18\pmod{27}$, not a square.  If $3\mid m$, then $3\nmid n$
and $F(m,n)$ is congruent to one of $8,17,26$, again not a square.

Finally suppose $3\nmid mn$.  Divide \eqref{eq:Fmod27} by the square
$n^6$ and put
\[
t\equiv(mn^{-1})^3\pmod{27}.
\]
The unit cubes modulo $27$ are
\[
1,8,10,17,19,26.
\]
For these six values,
\[
18t^2+6t+8\equiv5\ \text{or}\ 20\pmod{27},
\]
neither of which is a square.  Thus $F(m,n)$ is never a square.  Since
$x^2=\sigma F(m,n)$, the sign must be $\sigma=-1$.
\end{proof}

By \cref{prop:genus2}, the primitive projective solutions of
\[
x^2=-F(m,n)
\]
are
\[
[m:n:x]=[1:2:\pm4],
\quad
[2:3:\pm3],
\quad
[2:7:\pm895].
\]
The first has $n$ even and was excluded by \cref{lem:neven}.

For $[m:n]=[2:3]$, one has $3m-2n=0$.  Formula
\eqref{eq:paramsecond}, with the required norm sign, gives
$\delta=\alpha^2$ up to sign.  Therefore
\[
z+y\alpha
 =\etaunit\alpha^4
 =(\alpha^2-2)(3\alpha)
 =9-6\alpha.
\]
Thus
\[
(x,y,z)=(\pm3,-6,9),
\]
which is not primitive.

For $[m:n]=[2:7]$, the unscaled triple is
\[
(-112,-32,-41).
\]
With the required sign, we may take
\[
\delta=112+32\alpha+41\alpha^2.
\]
There is a useful internal square identity:
\[
\delta=(10+\alpha+2\alpha^2)^2.
\]
Indeed,
\[
(10+\alpha+2\alpha^2)^2
 =112+32\alpha+41\alpha^2.
\]
Moreover,
\[
\Norm(10+\alpha+2\alpha^2)
 =10^3+3+9\cdot2^3-9\cdot10\cdot1\cdot2
 =895.
\]
Hence $\Norm(\delta)=895^2$, so $x=\pm895$.

A further direct expansion gives
\[
\delta^2=20416+12211\alpha+10208\alpha^2.
\]
Multiplying by $\etaunit=\alpha^2-2$ yields
\begin{align*}
\etaunit\delta^2
&=(3\cdot12211-2\cdot20416)
 +(3\cdot10208-2\cdot12211)\alpha\\
&=-4199+6202\alpha.
\end{align*}
Therefore
\[
(y,z)=(6202,-4199),
\]
and the second unit class contributes exactly
\[
(\pm895,6202,-4199).
\]

\section{Completion of the proof}

By \cref{prop:squareclasses}, every primitive solution belongs to exactly
one of the two unit square classes $e=0$ or $e=1$.  The first class gives
only $(\pm1,0,1)$ after the nonprimitive candidate $(\pm3,3,0)$ is
removed.  The second gives only $(\pm895,6202,-4199)$ after the
nonprimitive candidate $(\pm3,-6,9)$ is removed.  This proves
\cref{thm:main}.

For completeness, the final direct checks are
\[
3\cdot0^3+1^3=1=(\pm1)^4
\]
and
\[
3\cdot6202^3+(-4199)^3
 =641641050625
 =895^4.
\]
Also
\[
\gcd(895,6202)=1,
\qquad
\gcd(895,4199)=1,
\qquad
\gcd(6202,4199)=1,
\]
so the large triple is primitive.

\appendix

\section{Finite congruence tables}

This appendix records the small residue computations used in the proof.

\subsection{Squares modulo 16, 27, and 64}

The square residues modulo $16$ are
\[
0,1,4,9.
\]
The square residues modulo $27$ are
\[
0,1,4,7,9,10,13,16,19,22,25.
\]
A square divisible by $16$ is congruent to $0$ or $16$ modulo $64$.

\subsection{Unit cubes modulo 27}

For $u\in(\Z/27\Z)^\times$, the possible values of $u^3$ are
\[
1,8,10,17,19,26.
\]
The expression used in \cref{lem:Fsign} has values
\[
\begin{array}{c|rrrrrr}
t&1&8&10&17&19&26\\ \hline
18t^2+6t+8\pmod{27}&5&20&5&20&5&20.
\end{array}
\]

\section{Independent exact checks}

The following identities can be verified entirely by expansion:
\begin{align*}
\operatorname{disc}(X^6+60X^3-72)
 &=2^{18}3^{25},\\
\operatorname{disc}(-8X^6-60X^3+9)
 &=2^{18}3^{25},\\
\operatorname{disc}(-F(X,1))
 &=2^{18}3^{25},\\
F(-U-V,-2U-V)&=-H(U,V),\\
H(T,1)-(4T^3-3T^2-24T+20)^2
 &=81(T-1)^3(3T+5).
\end{align*}
The finite-field counts in \cref{prop:genus2} are also obtained by direct
enumeration.  A companion Python file supplied with the manuscript performs
these checks using exact integer and finite-field arithmetic.

\section{Logical dependency summary}

For clarity, the classification rests on the following chain.

\begin{enumerate}
\item Primitivity implies pairwise coprimality and excludes the prime $3$
from $xz$.
\item The field $K=\Q(\sqrt[3]{3})$ has integral basis
$1,\alpha,\alpha^2$, class number $1$, and norm-$1$ units
$\etaunit^n$.
\item The principal ideal $(z+y\alpha)$ is a fourth power, so
$z+y\alpha=\etaunit^e\delta^2$ with $e\in\{0,1\}$ and
$\Norm(\delta)=x^2$.
\item The two values of $e$ produce two explicit conics in the
coefficients of $\delta$.
\item Primitive parametrization and local congruences reduce the problem
to rational points on $C_A,C_B,C_F$.
\item The exact genus-two point classification in \cref{prop:genus2}
leaves only the four primitive triples stated in \cref{thm:main}.
\end{enumerate}

\end{document}
